The general solvability assertion for nonsingular systems is commonly
called Howie's conjecture; see \cite[Introduction]{KlyachkoMikheenkoRomankov2024}
for this formulation and terminology. Its one-equation, one-variable
case is the Kervaire--Laudenbach conjecture: a nonzero exponent sum
should guarantee coefficient injectivity. This is stronger than the
nontriviality assertion usually called the Kervaire conjecture, which
asks whether adjoining one generator and imposing one relation can
annihilate a nontrivial group. Chen explains these formulations and their
connection with high-dimensional knot groups
\cite[Section~1]{Chen2026}. Theorem~\ref{thm:main} establishes the
finite nonsingular-system assertion for arbitrary coefficient groups.

The foundational system theorem is due to Gerstenhaber and Rothaus
\cite{GerstenhaberRothaus1962}. For a square system over a compact
connected Lie group, they compute the degree of the word map from the
determinant of its exponent-sum matrix. A nonzero determinant therefore
forces surjectivity and a simultaneous solution in that Lie group.
Their finite-group theorem also produces a finite solution overgroup,
using arithmetic specialization and reduction over finite fields.
A nonsingular rectangular
system reduces to the square case by retaining columns of a nonzero
maximal minor and setting the other variables equal to the identity. Their
argument established the usefulness of compact Lie groups, cohomology,
and degree theory in a problem stated entirely in group-theoretic terms.

A different line of work uses the topology of relative presentations.
Howie proved solvability of finite independent systems over locally
indicable groups \cite[Corollary~4.2]{Howie1981}; locally indicable
means that every nontrivial finitely generated subgroup maps onto
$\mathbb Z$. Klyachko's theorem treats a single equation in one variable
over a torsion-free group when the exponent sum is $\pm1$
\cite{Klyachko1993}. This unimodular hypothesis is stronger than a
nonzero exponent sum. Chen later recovered that theorem by estimating
the complexity of surfaces in HNN extensions
\cite[Theorem~6.9]{Chen2026}. These results illustrate two complementary
sources of control: algebraic conditions on the coefficient group and
topological restrictions on diagrams witnessing a kernel element.

The compact-unitary method also passes to metric ultraproducts. Pestov
observed its application to hyperlinear groups
\cite[Corollary~10.4]{Pestov2008}, and Nitsche and Thom give the
nonsingular-system statement explicitly
\cite[Theorem~1.2 and Lemma~2.1]{NitscheThom2022}. Here a hyperlinear,
or Connes-embeddable, group is one that embeds into a metric ultraproduct
of unitary groups with their normalized Hilbert--Schmidt metrics.
Their stronger Theorem~1.3 replaces nonsingularity by vanishing second
homology of a covering of the presentation complex obtained after
deleting the coefficients. For the presentation complex itself,
vanishing second homology is precisely independence of the exponent
rows. Their covering criterion therefore reaches beyond nonsingular
systems while retaining the hypothesis on the coefficient group.

Other developments control the solution overgroup more closely.
Klyachko, Mikheenko, and Roman'kov obtain solutions within specified
classes of solvable groups with suitable torsion-free abelian factors
\cite{KlyachkoMikheenkoRomankov2024}. Ramirez-C\^ot\'e and Wise use
the Banach fixed-point theorem for groups embedded in Magnus-type
power-series groups; their nonsingular-system construction inverts the
determinant in the coefficient ring
\cite[Theorems~9--10]{RamirezCoteWise2025}. Such conclusions address
additional structure beyond existence in an arbitrary overgroup.

The direct antecedent of our proof is the spectral-phase and planar
argument in \emph{The Kervaire theorem for groups}
\cite[Lemma~2.1, Lemma~3.1, and Theorem~4.1]{OpenAIKervaire2026}.
Its fixed-plane incidence construction treats one unimodular relator.
The extension below uses products of unitary groups and rational
intersection classes to accommodate several relator types with
different multiplicities. The cohomological background is classical
\cite{Borel1953}; related geometry of unitary eigenvalue strata and
their intersection classes is developed by Nicolaescu
\cite{Nicolaescu2010}. In finite dimensions, spectral-phase
subadditivity also follows from Thompson's exponential formula
\cite{Thompson1986}. Antezana, Larotonda, and Varela prove an
approximate version, in operator norm, for embeddable $\mathrm{II}_1$ factors
\cite[Theorem~4.4]{AntezanaLarotondaVarela2012}. Our argument instead
uses the faithful trace of the group von Neumann algebra
\cite[Section~3.3]{Jones2009} and proves the required phase inequality
directly, so the coefficients need no unitary approximation hypothesis.
